%!TEX root = cv.tex

\cvsection{Research Highlights}
\note{Five papers that best represent the contribution; the complete list
follows. \lead\ denotes lead or corresponding authorship.}

\vspace{4pt}
\begin{enumerate}

\item \textbf{GW170817: constraining the nuclear matter equation of state from
the neutron star tidal deformability} \quad \citecount{298}\\
T. Malik, N. Alam, M. Fortin, C. Provid\^encia, B. K. Agrawal, T. K. Jha,
\me, S. K. Patra \\
\jref{Phys.\ Rev.\ C 98, 035804 (2018)}\quad \pdoi{10.1103/PhysRevC.98.035804}\\
{\small Established the correlations that turn a measured tidal deformability
into bounds on nuclear incompressibility and symmetry energy --- one of the most
cited routes from GW170817 to laboratory nuclear physics.}

\item \lead\textbf{New relativistic effective interaction for finite nuclei,
infinite nuclear matter and neutron stars} \quad \citecount{153}\\
\me, B. K. Agrawal, S. K. Patra \\
\jref{Phys.\ Rev.\ C 97, 045806 (2018)}\quad \pdoi{10.1103/PhysRevC.97.045806}\\
{\small Introduced the IOPB-I interaction, now used internationally as a
standard relativistic mean-field parameterisation spanning finite nuclei to
neutron-star matter.}

\item \textbf{Effects of dark matter on the nuclear and neutron star matter}
\quad \citecount{116}\\
H. C. Das, A. Kumar, \me, S. K. Biswal, T. Nakatsukasa, A. Li, S. K. Patra \\
\jref{Mon.\ Not.\ R. Astron.\ Soc.\ 495, 4893 (2020)}\quad
\pdoi{10.1093/mnras/staa1435}\\
{\small First systematic treatment of WIMP dark matter inside the relativistic
mean-field description of dense matter; opened the group's ongoing programme on
dark-matter-admixed compact objects.}

\item \lead\textbf{Inferring neutron star properties from GW170817 with
universal relations} \quad \citecount{110}\\
\me, P. Landry \\
\jref{Phys.\ Rev.\ D 99, 123026 (2019)}\quad \pdoi{10.1103/PhysRevD.99.123026}\\
{\small Used equation-of-state-insensitive relations to extract radii, moments
of inertia and binding energies of the GW170817 components without assuming a
specific equation of state.}

\item \lead\textbf{Constraints on the moment of inertia of PSR
J0737$-$3039A from GW170817} \quad \citecount{93}\\
P. Landry, \me \\
\jref{Astrophys.\ J. Lett.\ 868, L22 (2018)}\quad
\pdoi{10.3847/2041-8213/aaee76}\\
{\small Predicted the moment of inertia of the double-pulsar primary ahead of
its measurement, linking gravitational-wave and radio-timing observables.}

\end{enumerate}
